\chapter{Steady-State System Design and Theoretical Framework\label{cha:chapter5}}

This chapter contains the theoretical framework and steady-state design for the autonomy stack of the warehousing fleet under nominal conditions. The design incorporates physical domain constraints and leverages edge processing for the runtime cognitive processing pipeline. This baseline consists of three interlocked levels. The first deals with energy-aware kinematics, balancing the robot's power consumption while relying on a regulated tracking paradigm \cite{macenski2020marathon}. The second implements decentralized coordination, enabling collision-free spatial navigation via peer-to-peer decentralized spatial mutexes and removing the single point of failure embodied by the centralized dispatchers \cite{ricart1981optimal}. The third level utilizes edge processing to enable real-time simulation environment parsing and the fused computer vision pipeline using uniform LBP and HOG descriptors for the real-time perception of the environment \cite{ojala2002multiresolution}. These three levels form a set of mutually dependent steady-state models, defining the high-level constraints for the multi-agent system's behavior. This operational envelope, therefore, serves as a baseline for the hierarchical fault-tolerant framework introduced in Chapter 6.

This particular structural workflow has been captured in the form of a steady-state system architecture pipeline, namely, Figure~\ref{fig:steady_state_architecture}. In this figure, nominal processing loops have been depicted, illustrating how low-level telemetry feeds into high-level distributed coordination and computer vision layers, creating a predictive, multi-tiered pipeline in which all stages are deterministic, i.e., the processing in one layer serves as an input boundary condition for processing in the subsequent layer.

\section{Hybrid Battery Modeling and Energy Consumption Formulations}

\subsection{Theoretical Foundations and Power Loss Trade-Offs}

A monolithic discharge curve that linearly depends on time spent in a mission confuses the difference between a robot that has been waiting at a waypoint and one that has been driving at full throttle down a narrow aisle; the only difference between the two is the amount of time spent performing the task. This does not account for the amount of torque applied and the resulting current draw by the wheels during rotations, especially during the tight turns in narrow aisle intersections mandated by the problem geometry. A robot that is parked in a vertical scanning state consumes power, but nowhere near the amount of current that the same robot would during the backward drive, rotation, or other maneuvers to escape a collision or inspection waypoints. Pooling the two discharge multipliers together loses the signal needed to accurately monitor energy depletion levels. A robot that spends most of a mission in idle inspection waypoints should not have a similar charge estimate to a robot that spends the same real time in dynamic translational and rotational trajectories

As a result, the hybrid discharge model separates the two power loss processes rather than considering a first-order approximation of the two as a single coefficient. First, a static term covers idle power consumption associated with non-translation motion as well as the background computational overhead of the onboard processing pipeline. Second, a term that depends on the distance driven captures the additional current draw by the drivetrain during rotations and forward motion. This choice reflects the telemetry resolution the battery monitoring system provides to the task allocation system on a $10\,\text{Hz}$ control loop rather than attempting to capture instantaneous current consumption as one would implement in a full electromechanical battery model as seen in \cite{mei2005case}.

\subsection{Mathematical Formalization of the Hybrid Decay Function}

To capture the physics of the discharge during regular warehouse operations, a hybrid decay function was implemented to emulate the battery's capacity as a function of the robot's energy expenditure throughout the mission cycle. In particular, the agent's energy consumption follows the non-linear hybrid decay model:
\begin{equation}
E(t) = E_{\text{start}} - \left( \lambda_{\text{idle}} \cdot t + \eta_{\text{dynamic}} \cdot \int_{0}^{t} \| \mathbf{v}(\tau) \|_2 \, d\tau \right) \label{eq:hybrid_decay_continuous}
\end{equation}
with the constant term accounting for the baseline power consumption by the on-board computer and the regular sensors suite while the tracking of the wheel rotation speed for the purpose of the real-time odometry introduces an unnecessary overhead to the edge computing firmware. With the selection of the simplified parametrization, the physical energy expenditure was then approximated with a reduced update rate proportional to the planar displacement between the successive AMCL filter updates:
\begin{equation}
\mathcal{T}_{\text{AMCL}} = \{ t_0, t_1, t_2, \dots, t_N \} \label{eq:amcl_time_set}
\end{equation}
Given the sequence of updates, the state vector of the system at the $k$-th update period $t_k \in \mathcal{T}_{\text{AMCL}}$ is represented as a planar position vector:
\begin{equation}
\mathbf{x}_k = [x_k, y_k]^T \in \mathbb{R}^2 \label{eq:amcl_pose_vector}
\end{equation}
with $x_k, y_k$ denoting the AMCL-estimated coordinates at update time $t_k$. Then, the planar distance between the successive updates can be computed as:
\begin{equation}
\Delta d_k = \| \mathbf{x}_k - \mathbf{x}_{k-1} \|_2 = \sqrt{(x_k - x_{k-1})^2 + (y_k - y_{k-1})^2} \label{eq:euclidean_displacement}
\end{equation}
via the substitution of the filtered position estimates for the raw wheel rotation speed values, effectively cutting out the high-frequency components from the position signal. With the definition of the spatial displacement between the AMCL filter updates, the hybrid discharge function can then be re-written in the discrete-time domain for the $N$-step mission cycle as:
\begin{equation}
E(t_N) = E_{\text{start}} - \left( \lambda_{\text{idle}} \cdot t_N + \eta_{\text{dynamic}} \cdot \sum_{k=1}^{N} \Delta d_k \right) \label{eq:hybrid_decay_discrete}
\end{equation}

In order to avoid artificial inflation of the cumulative mission energy, the state monitoring engine incorporates a maximum displacement threshold into consideration. In particular, one often observes during simulation runs that the initialization artifacts lead to unrealistically large position jumps, e.g., in AMCL. This effect, if unaccounted for, would artificially inflate the mileage information, resulting in the battery model cutting-off the agents' activity prematurely. This issue is addressed by limiting the step size in $\Delta d_k$ using a simple non-linear filter:
\begin{equation}
\Delta d_{k, \text{filtered}} = \begin{cases} \Delta d_k & \text{if } \Delta d_k < \delta_{\text{warp}} \\ 0 & \text{if } \Delta d_k \ge \delta_{\text{warp}} \end{cases} \label{eq:teleportation_filter}
\end{equation}
which ensures that the step size does not exceed $\delta_{\text{warp}} = 5.0\text{ m}$. This is a reasonable bound, given the differential constraints of the robot platform, which renders most of the non-physical position jumps in $\Delta d_k$ filtered out, thus preventing the artificial inflation of the cumulative mission energy in Eq.~\ref{eq:hybrid_decay_discrete}.

While this discrete-time representation mathematically characterizes the energy states at major decision epochs, the onboard battery monitoring module implements this formulation via decoupled tracking threads. Specifically, the static idle decay term is evaluated continuously on a fixed $10\,\text{Hz}$ local telemetry control loop, whereas the dynamic mileage integration accumulates discrete displacement steps between consecutive estimator publications, ensuring that micro-movements are continuously accumulated. For the particular choice of the robot platform, the value of $\lambda_{\text{idle}}$ was set to induce a linear discharge of $0.02\,\%\,\text{s}^{-1}$ during idle periods, corresponding to approximately 0.002\% capacity loss per $10\,\text{Hz}$ control loop iteration. Furthermore, the dynamic component of the hybrid discharge model was calibrated with the coefficient $\eta_{\text{dynamic}} = 0.15\,\%\,\text{m}^{-1}$ to account for the effective current draw characteristics of the drivetrain on the concrete floor of the warehouse facility. By virtue of the simplified parametrization, the computation of the hybrid discharge model is performed at the battery management level with minimal overhead, thus bypassing the substantial CPU and memory overhead of the raw telemetry parsing and message unpacking at the high update rates described in \cite{mei2005case}.

\section{Spatial Mutex Mechanics and Peer-to-Peer DDS Coordination}

\subsection{DDS Transport Contract Formulations and QoS Robustness}

To resolve this issue for multi-agent telemetry, the diagnostic framework has to establish a consistent communication contract with the distributed nodes. Ideally, the state estimation publications should be subscribable under varying network conditions while preserving enough computational resources for the passive observers. To address these constraints, the diagnostic framework implements a dual transport contract, which is produced redundantly by the offline diagnostic suite. While the active inference agent only subscribes to the high-level quality of service $\mathcal{Q}_{\text{TL}}$, the diagnostic suite maintains both volatile best-effort $\mathcal{Q}_{\text{BE}}$ and transient-local reliable $\mathcal{Q}_{\text{TL}}$ contracts. As a result, the AMCL pose publications can always be observed independently from the cold start QoS pinning of the active edge agents. By design, the diagnostic telemetry does not add to the computational burden of the active inference agents while providing sufficient reliability for the passive observers.

\subsection{Decentralized Coordination vs. Centralized Failure Modes}

Multi-agent fleet management typically relies on centralized dispatch servers that allocate trajectories and resolve spatial conflicts. While such approaches are typically computationally light, they have a single point of failure and require uninterrupted communications from every robot in the fleet to the master controller. In practice, this is a questionable assumption in large-scale warehouse environments where deep steel shelves create wireless dead zones and multi-path fading. A central controller cannot safely arbitrate collisions in a steel-enclosed aisle if one of the robots in it has lost radio contact. Our decentralized P2P coordination approach breaks this centralized single point of failure by pushing the arbitration logic into the agents. Acting as independent edge nodes, the robots coordinate their movements on the industrial floor utilizing a decentralized, brokerless data distribution paradigm native to the fleet's coordinated communication framework \cite{maruyama2016exploring}. 

\subsection{Corridor Spatial Mutex Formulations and Shared Corridors}

The central problem of a decentralized approach is preventing deadlocks in confined spaces. Whereas most open regions allow for sufficient lateral space for two robots to safely pass each other, narrow aisles have insufficient width for such a maneuver. Specifically, the narrow aisle width was found to be less than twice the robot's safety radius plus an epsilon-long robot-specific margin \cite{ricart1981optimal}. Therefore, the local costmaps of robots approaching each other from opposite ends of the aisle would erroneously inflate their respective walls and the other robot's body, making mutual passage impossible.

To formalize this into an algorithm, we define the narrow aisle spatial mutex at the corridor level, not the inspection region level. We identify three narrow shared corridors on the factory floor plan, all aligned parallel to the north-south axis and distributed along the east-west spatial span. Each narrow aisle is shared between two sets of shelves, spanning the positive and negative lateral faces of the adjacent large shelving units, respectively. To pass the shared narrow corridors, the robot must make a claim on the relevant shelf face, which then makes this whole aisle region mutually exclusive. By contrast, the virtual mobile cluster is located at the wide loading dock, where the clearance is sufficient for robot traffic from both sides of the aisle. Hence, the task allocator does not apply the shared corridor mutex to the mobile cluster.

To avoid spatial conflicts in the narrow part of the transit zone, the workspace is divided into three sections along the length of the corridor. The sections are named the western, middle, and eastern segments of narrow aisle areas. Each section is limited by pairs of facing shelf units whose approach from both sides would lead to local costmap inflation. The first segment, designated as the western corridor, is situated between the positive longitudinal face of the first large shelving unit and the negative longitudinal face of the adjacent fifth large shelving unit. The second segment, designated as the middle corridor, is located between the positive longitudinal face of the fifth large shelving unit and the negative longitudinal face of the fourth large shelving unit. The third segment, designated as the eastern corridor, is located between the positive longitudinal face of the fourth large shelving unit and the negative longitudinal face of the third large shelving unit. The coordination module restricts navigation with these pairs of shelf faces, ensuring exclusive access to the entire length of the narrow aisle corridor, avoiding deadlocks by means of distributed mutual exclusion \cite{ricart1981optimal}.

\subsection{Peer-to-Peer Claims Protocol and Lease Mechanics}

The coordination of the entry into narrow corridor segments follows the asynchronous distributed state broadcasts mechanism. Instead of having a single coordinator regulating the claims procedure, each agent individually keeps track of the relevant information. To check whether a target segment is available to the current evaluating agent, a Boolean product formula is used to determine the overall segment's availability for use. Given an arbitrary corridor segment, its availability for use by an evaluating agent at a particular moment is defined using the Boolean product formula:
\begin{equation}
\mathcal{A}_i(C, t) = \prod_{j \neq i} \left( 1 - \mathcal{C}_j(C, t) \cdot \mathbb{I}\left( t < t_{\text{expire}, j} \right) \right) \label{eq:spatial_availability_check}
\end{equation}
where the active claim of a peer is defined by:
\begin{equation}
\mathcal{C}_j(C, t) \in \{0, 1\} \label{eq:active_claim_status}
\end{equation}
with the indicator function being equal to unity when the lease of the particular peer is active, i.e., when the current moment is before the expiration deadline specified for the lease. The value of the deadline time itself is given by:
\begin{equation}
t_{\text{expire}, j} = t_{\text{stamp}, j} + \tau_{\text{lease}} \label{eq:lease_expiration_time}
\end{equation}
where the default value during normal operations is set to $\tau_{\text{lease}} = 6.0\,\text{s}$. Note that to sustain the lease while waiting for tracking and inspection procedures, an active claiming agent needs to renew the lease by sending updated lease information to the fleet's communication infrastructure. The parameter governing the time interval between the consecutive renewal messages is set to $T_{\text{renew}} = 1.5\,\text{s}$. The deadline time is moved to a new value specified by Equation \ref{eq:lease_expiration_time}, thus preventing the corridor segment from being released while the claiming agent sustains the lease. If the lease renewal procedure is not followed, the indicator function drops to zero, and the product formulation given by Equation \ref{eq:spatial_availability_check} evaluates to unity, automatically restoring the segment's availability for the fleet's other agents.

\subsection{Asymmetric Race-Condition Resolution}

To resolve simultaneous claim packets in the asynchronous distributed architecture, we define the tie-breaking priority relation:
\begin{equation}
(t_{\text{stamp}, i}, \; \text{ID}_i) \prec (t_{\text{stamp}, j}, \; \text{ID}_j) \iff \left( t_{\text{stamp}, i} < t_{\text{stamp}, j} \right) \lor \left( (t_{\text{stamp}, i} = t_{\text{stamp}, j}) \land (\text{ID}_i < \text{ID}_j) \right) \label{eq:asymmetric_precedence}
\end{equation}
which imposes a linear ordering of the packets by their arrival time and agent-specific identifier. Practically however, it is impossible to enforce such ordering on the edge due to serialization overhead on the network. The physical peer-to-peer stack uses a stateless local runtime coordination mechanism, which leverages LCFS (Last-Come-First-Served) lease preemption. By doing so, incoming lease updates can preemptively overwrite local lock registries, obviating the need to maintain a queue at the edge that would serialize updates on the planar manifold, conserving critical edge compute cycles and reducing coordination induced latency.

\section{Runtime SDF Parsing and 2D Coordinate Transformation Geometry}

\subsection{Dynamic Simulation Environment Parsing vs. Static Layouts}

The utilization of pre-set, static waypoint lists severely limits the versatility of multi-agent warehousing fleets. In a scenario of reconfiguration or re-positioning of the stored items, the list becomes essentially useless and has to be updated by the system administrators, which is a very time-consuming and error-prone task. Subsequently, wrong destination coordinates may lead to mission abortion due to local planning errors or collisions. The previously described limitations are addressed by the current solution, which implements a dynamic simulation environment parsing framework. Upon launch, the coordinate resolution engine instantiates the world parsing node, which parses the simulation world layout representation directly for the purposes of dynamic coordinate extraction. Using hierarchical document parsing techniques, the spatial parser node tracks down the first-level include tags, inside of which the necessary model frames are located. For each discovered frame, the spatial parser node extracts the raw, six-degrees-of-freedom pose vector relative to the global origin:

\[ \mathbf{P}_{\text{raw}} = \begin{bmatrix} x_{\text{raw}} & y_{\text{raw}} & z_{\text{raw}} & \phi_{\text{roll}} & \theta_{\text{pitch}} & \psi_{\text{yaw}} \end{bmatrix}^T \]

Given the structure of the raw simulation world layout file, the target model names are parsed according to the small shelf/big shelf/pallet pattern for dynamic coordinate replacement. Given that the multi-agent system operates purely within a planar SE(2) manifold within a plane defined by the $XY$ axes, the spatial parser node discards the vertical position vector component $z_{\text{raw}}$ as well as the roll and pitch angles of the frame, effectively obtaining the reduced $\text{SE}(2)$ representation of the target position:

\[ 
\mathbf{q}_{\text{projected}} = \begin{bmatrix} x_{\text{raw}} & y_{\text{raw}} & \psi_{\text{yaw}} \end{bmatrix}^T \in \text{SE}(2) 
\]

For the purposes of the current project, this reduced, planar representation of spatial coordinates is sufficient to implement a non-holonomic kinematic chain, thereby obtaining a dynamic path planning framework not limited by the static waypoint list.

\subsection{Asymmetric Model Origin Corrections and Homogeneous 2D Rotations}

The model coordinates specified directly from the layout simulation software represent the origin as a point in the uncalibrated CAD-centered reference frame. For physical shelving equipment, these default frames do not align with the actual geometric center of the three-dimensional structures they represent. To extrapolate true 2D positions, the coordinate resolution engine utilizes homogeneous coordinate transforms on the planar $\text{SE}(2)$ manifold. Let the global world frame be denoted as $\mathcal{G}$, the uncorrected default model frame be $\mathcal{O}$, and the corrected geometric center be $\mathcal{L}$. The position and orientation of the model origin frame $\mathcal{O}$ relative to the global map frame $\mathcal{G}$ is represented with the following homogeneous transform:
\begin{equation}
\mathbf{T}_{\mathcal{G}}^{\mathcal{O}} = \begin{bmatrix} \cos\psi_{\text{yaw}} & -\sin\psi_{\text{yaw}} & x_{\text{raw}} \\ \sin\psi_{\text{yaw}} & \cos\psi_{\text{yaw}} & y_{\text{raw}} \\ 0 & 0 & 1 \end{bmatrix} \label{eq:homogeneous_raw_transform}
\end{equation}
To encode the asymmetric offset from the CAD world origin to the actual center of the physical shelving model, a model-dependent local translation offset is defined in homogeneous coordinates as:
\begin{equation}
\mathbf{p}_{\text{local}} = \begin{bmatrix} \Delta x_{\text{local}} & \Delta y_{\text{local}} & 1 \end{bmatrix}^T \label{eq:local_offset_vector}
\end{equation}
where the local translations have been parameterized according to the layout specification. In particular, small and big shelving units are set to have a longitudinal offset of $-0.5\,\text{m}$ and lateral offset of $0.0\,\text{m}$, while static pallets do not require any additional transformation. Using the projection expressed by the homogeneous offset vector and transformation matrix, the true corrected geometric center vector relative to the global map frame $\mathcal{G}$ can be calculated as:
\begin{equation}
\mathbf{x}_{\mathcal{G}} = \begin{bmatrix} cx & cy & 1 \end{bmatrix}^T \label{eq:corrected_center_vector}
\end{equation}
where the projection is realized by the following matrix multiplication:
\begin{equation}
\mathbf{x}_{\mathcal{G}} = \mathbf{T}_{\mathcal{G}}^{\mathcal{O}} \mathbf{p}_{\text{local}} \label{eq:homogeneous_offset_projection}
\end{equation}
This effectively realizes a rotation of the local offset translation by the model's yaw angle and addition of the vector to the global raw coordinate, resulting in a coordinate anchor suitable for subsequent waypoint computations. Finally, to define the target inspection pose of the robot footprint frame $\mathcal{R}$ relative to the corrected geometric center frame $\mathcal{L}$, the relative homogeneous transformation is defined as:
\begin{equation}
\mathbf{T}_{\mathcal{L}}^{\mathcal{R}} = \begin{bmatrix} \cos\psi_{\text{rel}} & -\sin\psi_{\text{rel}} & \Delta x_{\text{approach}} \\ \sin\psi_{\text{rel}} & \cos\psi_{\text{rel}} & \Delta y_{\text{approach}} \\ 0 & 0 & 1 \end{bmatrix} \label{eq:homogeneous_relative_transform}
\end{equation}
which enables the computation of the compound homogeneous transform representing the robot's final pose in the global map frame $\mathcal{G}$ as:
\begin{equation}
\mathbf{T}_{\mathcal{G}}^{\mathcal{R}} = \mathbf{T}_{\mathcal{G}}^{\mathcal{L}} \mathbf{T}_{\mathcal{L}}^{\mathcal{R}} \label{eq:homogeneous_matrix_chain}
\end{equation}
By formulating the homogeneous transformations for local offset, rotation, and robot-specific transformation in this manner, the spatial parser node guarantees kinematic compatibility of the subsequent approach pose and camera transform calculations relative to the non-holonomic constraints of the agents.

\subsection{Parametric Approach Pose Generation and Steel-Column Structural Bypass}

To allow the vehicle's visual perception system to capture the necessary visual data, the robot must approach the target shelf at a sufficient distance $D_{\text{approach}} = 1.20\,\text{m}$ normal to its longitudinal inspection face. For small shelf instances (horizontal sweeping), this offset is determined by the model's depth parameter, yielding $D_{\text{offset}} = \frac{\text{depth}}{2} + D_{\text{approach}} = \frac{0.6}{2} + 1.2 = 1.50\,\text{m}$. The inspection waypoints are distributed along the longitudinal axis, with local division steps defined as $\Delta x_{\text{local, i}} = -1.35 + (i \cdot 0.9)$ for $i \in \{0, 1, 2, 3\}$. The global inspection target coordinate vector $\mathbf{q}_{\text{small}, i} = [gx_{s, i}, gy_{s, i}]^T$ is resolved by applying the rigid-body transformation:

\[ 
\begin{bmatrix} gx_{s, i} \\ gy_{s, i} \end{bmatrix} = \begin{bmatrix} cx \\ cy \end{bmatrix} + \begin{bmatrix} \cos(\psi_{\text{yaw}}) & -\sin(\psi_{\text{yaw}}) \\ \sin(\psi_{\text{yaw}}) & \cos(\psi_{\text{yaw}}) \end{bmatrix} \begin{bmatrix} \Delta x_{\text{local, i}} \\ \delta_s \cdot D_{\text{offset}} \end{bmatrix} 
\]

where $\delta_s \in \{-1, 1\}$ represents the directional approach sign designating either the positive or negative lateral side of the shelving unit.

Conversely, the big shelf type (vertical sweeping) utilizes a lateral axis projection where the approach offset is determined by the model's width parameter, giving $D_{\text{offset}} = \frac{\text{width}}{2} + D_{\text{approach}} = \frac{2.1}{2} + 1.2 = 2.25\,\text{m}$. The longitudinal section dividers are placed at $\Delta y_{\text{local, i}} = 6.75 - (i \cdot 4.5)$ relative to the model's origin. The global approach pose vector $\mathbf{q}_{\text{big}, i} = [gx_{b, i}, gy_{b, i}]^T$ for each inspection section is computed via:

\[ 
\begin{bmatrix} gx_{b, i} \\ gy_{b, i} \end{bmatrix} = \begin{bmatrix} cx \\ cy \end{bmatrix} + \begin{bmatrix} \cos(\psi_{\text{yaw}}) & -\sin(\psi_{\text{yaw}}) \\ \sin(\psi_{\text{yaw}}) & \cos(\psi_{\text{yaw}}) \end{bmatrix} \begin{bmatrix} \delta_s \cdot D_{\text{offset}} \\ \Delta y_{\text{local, i}} \end{bmatrix} 
\]

By mathematically separating these sweeping geometries, the coordinate transformation engine guarantees precise, collision-free waypoint generation that respects the non-holonomic kinematics of the fleet.

Nevertheless, due to the presence of permanent objects in the warehouse, some of the waypoints calculated using the parametric method might have intersected with the obstacles' inflation layers. An illustrative example is the steel columns presented at the $Y$-axis boundary of the facility, namely:
\begin{equation}
Y \approx -15.00\,\text{m} \label{eq:column_obstacle_pos}
\end{equation}
Indeed, according to the nominal coordinates, the approach section waypoint of the third inspection zone, namely, section C, for both the positive longitudinal face of the first large shelving unit and the negative longitudinal face of the fifth large shelving unit, appeared within the inflation boundary of the indicated columns. As a result, the global path planner was unable to find a clear path from the robot's start position to the generated approach waypoint due to the local costmap inflation. To resolve the navigation conflict, the spatial coordinate generator utilized a heuristic mapping rule to safely shift the target pose outside the inflation boundaries of the local costmap. In particular, for the sections of large shelves presented on the positive longitudinal face of the first large shelving unit and the negative longitudinal face of the fifth large shelving unit, when generating the approach waypoints to their section C, the global coordinates of the robot were shifted along the parallel axis to a safe distance:
\begin{equation}
Y_{\text{shifted}} = -13.50\,\text{m} \label{eq:column_shifted_pos}
\end{equation}
from the steel columns for visual inspection. The waypoint deviation ensured a collision-free inspection while maintaining the robot's visual direction normal to the large shelf's surface.

Although most of the inventory units allowed the symmetric approach to the inspection section from both sides, several storage units located near the permanent infrastructure required a specific approach displacement. In particular, in the example presented by the designated near-boundary pallet station, a local inventory unit located immediately next to the permanent wall could not utilize the standard approach offset:
\begin{equation}
1.81\,\text{m} \label{eq:standard_pallet_offset}
\end{equation}
from the positive side. Indeed, such a displacement would require the robot to finalize its approach to the inspection section within the inflated collision area of the warehouse's wall. To ensure a collision-free approach to the inventory unit while maintaining the robot's viewpoint normal to the inspection section, the coordinate estimator utilized a modified inspection approach. In particular, while the robot maintained the standard approach offset of $1.81\,\text{m}$ from the negative side, the positive-side approach was redirected with an orthogonal standing offset of: 1.60 m. To maintain strict alignment with the physical configuration space of the fleet, this displacement is calculated as a sum of the half pallet width of 0.40 m (derived from the full model width of 0.80 m defined in the configuration specs) and the nominal approach clearance of 1.20 m, yielding exactly the required 1.60 m clearance. As a result, the approach waypoint provided an unobstructed inspection vector $(x_{\text{approach}}, y_{\text{approach}})$ with an orthogonal viewpoint for the visual inventory check.

\subsection{Dynamic Mobile Clustering Formulations}

We face a similar spatial allocation problem while determining approach poses for the dynamic inventory items scattered across the loading dock. Decomposing the task into individual mobile pallets will cause redundancy in terms of navigation waypoints and collision risks in the crowded workspace. In order to solve this problem, the coordinate resolution engine creates a mixed cluster reference frame that tracks the mobile inventory's centroid:

\begin{equation} 
    \mathbf{c}_{\text{cluster}} = \left( \frac{1}{M} \sum_{k=1}^{M} x_k, \quad \frac{1}{M} \sum_{k=1}^{M} y_k \right) \label{eq:cluster_centroid} 
\end{equation}

where $M$ stands for the number of mobile pallets detected in the workspace, and $(x_k, y_k)$ represents the center coordinates of the $k$-th pallet. The list of inventory items is then sorted by the global $Y$-axis to impose a sequential north-south navigation pattern. To enforce consistent orientation constraints on the local path tracking controller, we enforce the aligned zero-yaw approach to all mobile pallets:
\begin{equation} 
    \psi_{\text{yaw}} = 0.0\,\text{rad} \label{eq:zero_yaw_constraint} 
\end{equation}

Using this simplification, the rotation-free projection becomes possible, effectively reducing $\text{SE}(2)$ transformations to a single $X$-axis offset:

\begin{equation}
    x_{\text{waypoint}, k} = x_k + \delta_s \cdot D_{\text{approach}} \label{eq:cluster_x_offset} 
\end{equation}

\begin{equation} 
    y_{\text{waypoint}, k} = y_k \label{eq:cluster_y_offset} 
\end{equation}

where $\delta_s \in \{-1, 1\}$ becomes a directional sign coefficient depending on the target face, while $D_{\text{approach}} = 1.2\,\text{m}$ defines the standard approach distance in a perpendicular direction. The simplified projection follows the same pattern as the shelf inventory, creating a sweeping lane coordinate frame. Notably, the dynamic inventory approach poses are removed from the global costmap queue once resolved. This step is essential to prevent the global path planning algorithm from treating active inspection targets as impassable barriers

\subsection{Orientation-Driven Sweeping for Dead-End Wall Mitigation}

Calculating the sectional approach waypoints in a parametric manner can cause severe motion planning issues near the physical boundaries of the simulation domain. For example, the concrete walls on the eastern ($X \ge 10.0\,\text{m}$) and northern ($Y \ge 15.0\,\text{m}$) edges of the environment restrict the robot's movements, forcing it to perform complex rotational movements when completing the last sweep section in order to retreat from proximity to the walls.

The local planner resolves such issues by utilizing an orientation-driven physics-informed sweeping rule, which adds no computational overhead to the global planning stage. For example, if a target shelf is located near the eastern wall and satisfies the spatial boundary condition ($\bar{x}_{\text{model}} > 10.0\,\text{m}$), the scanning order is reversed depending on which face of the target the robot is approaching from. If the target is being approached from the north, the sweep will follow a west-facing pattern \cite{macenski2020marathon}:
\[ D \rightarrow C \rightarrow B \rightarrow A \]
while approaching the southern side of the same shelving unit will produce an east-facing scanning pattern:
\[ A \rightarrow B \rightarrow C \rightarrow D \]
A similar scanning pattern is applied to shelves along the northern boundary satisfying the northern boundary constraint ($\bar{y}_{\text{model}} > 15.0\,\text{m}$):
\[ D \rightarrow C \rightarrow B \rightarrow A \quad (\text{Southward sweep}) \]

The target sequence is reversed before being fed into the navigation framework to allow the robot to avoid the most constrained areas of the environment first. By prioritizing these most constrained shelves, the next set of navigation waypoints are gradually positioned closer to central, unconstrained locations. This orientation-driven sweeping heuristic is designed to mitigate motion oscillations and global planning failures in highly constrained dead-end areas, the closed-loop performance and quantitative benefits of which are experimentally evaluated and benchmarked in Section 8.3.

\section{Path Planning and Regulated Pure Pursuit Control Theory}

\subsection{Non-Holonomic Kinematic Constraints and Chassis Jacobian}

To make sure that the motion planning algorithms are generating parametrized trajectories that can actually be physically tracked by the fleet agents, it becomes necessary to bake in the constraints imposed by the differential chassis of the robots. The robotic footprint is modeled as a non-holonomic rigid body on a rigid floor plane corresponding to the warehouse floor while the no-lateral-slip condition holds true for wheel-ground contact patches. The no-lateral-slip condition, which prevents translation of the chassis perpendicular to the robot's longitudinal axis, is expressed as:
\begin{equation}
\dot{x}\sin\psi - \dot{y}\cos\psi = 0 \label{eq:non_holonomic_constraint}
\end{equation}
where the derivatives of $x$ and $y$, and the $\psi$ term represent the spatial velocity coordinates and angular position of the robotic footprint frame relative to the map frame. The forward-kinematics expressions of the differential chassis can then be used to relate the instantaneous task-space velocity of the robotic footprint to the angular velocities of the right and left wheels as:
\begin{equation}
\begin{bmatrix} \dot{x} \\ \dot{y} \\ \dot{\psi} \end{bmatrix} = \begin{bmatrix} \frac{r}{2}\cos\psi & \frac{r}{2}\cos\psi \\ \frac{r}{2}\sin\psi & \frac{r}{2}\sin\psi \\ \frac{r}{b} & -\frac{r}{b} \end{bmatrix} \begin{bmatrix} \dot{\theta}_R \\ \dot{\theta}_L \end{bmatrix} \label{eq:forward_kinematics_matrix}
\end{equation}
where $r = 0.033\,\text{m}$ is the nominal wheel radius while $b = 0.287\,\text{m}$ corresponds to the nominal wheel track width. The angular velocities of the right and left wheels are expressed as $\dot{\theta}_R$ and $\dot{\theta}_L$. This yields the Jacobian for the kinematics chassis as:
\begin{equation}
\mathbf{J}(\psi) = \begin{bmatrix} \frac{r}{2}\cos\psi & \frac{r}{2}\cos\psi \\ \frac{r}{2}\sin\psi & \frac{r}{2}\sin\psi \\ \frac{r}{b} & -\frac{r}{b} \end{bmatrix} \label{eq:chassis_jacobian}
\end{equation}
which takes the form of a linear operator between the joint velocity space and task-space velocity manifold of the planar $\text{SE}(2)$ as expressed by:
\begin{equation}
\dot{\mathbf{q}} = \mathbf{J}(\psi) \mathbf{u} \label{eq:jacobian_velocity_mapping}
\end{equation}
where $\mathbf{u} = [\dot{\theta}_R, \dot{\theta}_L]^T$. The above mapping allows the lookahead controllers to respect the robot's physical joint limits and velocity constraints while generating the required reference angular velocities for the wheels.

\subsection{Multi-Layer Costmap Hierarchies and Global Planning}

Autonomous navigation in dense, warehouse settings is well-suited to a decoupling strategy between global path-planning and local obstacle avoidance, achieved by maintaining a hierarchy of costmap layers \cite{macenski2020marathon}. A global planner and local controller run in parallel, with the former computing a trajectory in the map frame to a user-defined approach pose using a relatively large inflation radius $r_{\text{global\_inflation}} = 0.55\,\text{m}$ and low decay cost scaling to motivate centerline preference in broad, open regions. In particular, the global planner is invoked at the time of a target approach request, and selects a shortest-path heuristic grid-based search before engaging the local tracking loop.

By contrast, the approach to dynamic and static shelves for inspection requires the robot to enter at close ranges, making global planner centerline preference incompatible with the required maneuvering. This necessitates an aggressive inflation threshold in the local costmap frame to avoid obstacles while maneuvering in proximity to shelves. The local costmap layer is centered in the robot frame with an inflation radius of $r_{\text{local\_inflation}} = 0.28\,\text{m}$, which is $1\,\text{cm}$ greater than the robot's bounding radius of $r_{\text{robot}} = 0.27\,\text{m}$: a minimally-safe, paper-thin margin representing a $1\,\text{cm}$ clearance. This creates undesirable tracking difficulties for the local planner, resulting in frequent violations of the local costmap constraints, with the non-holonomic chassis experiencing frequent large-scale collision warnings and control loop oscillations.

\subsection{Adaptive Lookahead Dynamics and Yaw Regulation in Regulated Pure Pursuit}

The local path-tracking loop must conform the robot pose to the global plan with high fidelity, but short-distance pure pursuit lookahead regulation results in excessively stiff rotational responses in confined sections. By contrast, long-distance lookahead tracking results in the robot skipping through narrow sections of the path, cutting corners and incurring collision risk with shelf edges. To balance these trade-offs dynamically, the regulated pure pursuit controller continuously scales the lookahead distance as a function of the instantaneous forward velocity of the non-holonomic chassis. This adaptive lookahead distance is formulated as:
\begin{equation}
l(v) = \max\left(l_{\text{min}}, \; \min\left(l_{\text{max}}, \; v \cdot t_{\text{lookahead}}\right)\right) \label{eq:adaptive_lookahead}
\end{equation}
where the lookahead time gain is parameterized to $t_{\text{lookahead}} = 1.5\,\text{s}$, and the lookahead bounds are strictly constrained by the minimum lookahead limit of $l_{\text{min}} = 0.35\,\text{m}$ and the maximum lookahead limit of $l_{\text{max}} = 1.20\,\text{m}$. To prevent the tight 1cm local costmap margin from triggering collision recovery sequences during in-place rotations, the translation and rotation controllers are decoupled during large-yaw maneuvers. The translational motion of the robot is suspended when the heading error of the robot with respect to the lookahead point exceeds a threshold of 0.785 rad, forcing an in-place rotation with an angular acceleration bounded to 1.5 rad/s$^2$ to prevent excessive rotation and wheel slip.

\subsection{Curvature and Cost-Based Velocity Regulation}

When navigating through tight spaces at very high speeds, skidding might cause an unrepresentative kinematic velocity estimate in the differential drive odometry. The following strategy aims to reduce velocity in accordance with sharp turns in order to avoid skidding. In accordance with the curvature-regulated path tracking model, the reference linear velocity is scaled as:
\begin{equation}
v_{\text{scaled\_curvature}} = \max\left(v_{\text{min\_regulated}}, \; v_{\text{max}} \cdot \frac{r_{\text{curr}}}{r_{\text{min\_radius}}}\right) \label{eq:nav2_rpp_curvature_scaling}
\end{equation}

where $v_{\text{min\_regulated}} = 0.15\,\text{m/s}$ represents the regulated scaling velocity floor and \\ $r_{\text{min\_radius}} = 0.75\,\text{m}$ defines the minimum curvature radius threshold below which velocity scaling activates. When the instantaneous radius of curvature $r_{\text{curr}} = 1/\kappa$ falls below this threshold, the linear velocity is scaled down to mitigate wheel slip and prevent unrepresentative velocity estimations in differential drive odometry. Additionally, the scaling dynamics incorporate a curvature regulation gain constant $K_{\text{reg}} = 2.5$ to govern the rate of deceleration under non-linear constraints.

To balance path-tracking stability and task completion, the tracking controller maintains two distinct velocity floors. The curvature and local costmap regulation floor \small{($v_{\text{min\_regulated}}$)} is constrained to $0.15\,\text{m/s}$ to ensure safe cornering. Conversely, the absolute minimum translational approach floor ($v_{\text{min\_approach}}$) is set to $0.03\,\text{m/s}$. This lower threshold allows the chassis to execute slow, precise docking adjustments near shelving units during the terminal waypoint alignment phase without triggering false-positive stall behaviors in the local planner.

In order to avoid collisions with objects along the course due to potential lateral drift, the forward velocity is additionally scaled down according to local costmap information:
\begin{equation}
v_{\text{scaled\_cost}} = v_{\text{scaled}} \cdot \left(1.0 - \frac{C}{C_{\text{max}}}\right)^{\gamma} \label{eq:cost_scaling}
\end{equation}
with $C$ being the cost of the cell of the costmap where the robot's footprint is currently located on, 254 being the maximum lethal obstacle cost value, and the scaling exponent $\gamma$ being set to a unit value of 1.0.

\subsection{Odometry Covariance Calibration and State Estimation}

To prevent wheel odometry from overwhelming the state estimator with unreasonably low variance values, the raw covariance estimates are calibrated. First, mathematically overconfident, zero-noise wheel odometry covariance matrices from standard kinematic models are discarded. Then, a calibrated diagonal covariance matrix that takes into account local skidding on concrete floors is applied to the processed odometry estimates:
\begin{equation}
\boldsymbol{\Sigma}_{\text{calibrated}} = \text{diag}\left(4\cdot10^{-4},\; 4\cdot10^{-4},\; 4\cdot10^{-4},\; 1\cdot10^{-4},\; 1\cdot10^{-4},\; 1\cdot10^{-2}\right) \label{eq:calibrated_covariance_matrix}
\end{equation}

In order to have a physical correspondence, the diagonal elements of the calibrated covariance matrix were associated with the six-dimensional state vector $\mathbf{x}_{\text{state}} = [x, \; y, \; z, \; \phi, \; \theta, \; \psi]^T \in \mathbb{R}^6$. The first three diagonal elements, each scaled to $4\cdot10^{-4}\,\text{m}^2$, define the translational variances along the three coordinate axes. The fourth and fifth diagonal elements were each set to $1\cdot10^{-4}\,\text{rad}^2$ and define the out-of-plane rotational variances of the roll and pitch states. The motivation for minimizing the out-of-plane elements was to prevent large floor discontinuities from artificially biasing the flat plane state. The sixth diagonal element was set to $1\cdot10^{-2}\,\text{rad}^2$ and defined the planar yaw heading variance that was temporarily inflated to allow wheel slip during hard turns without inducing filter divergence.

The calibrated diagonal covariance matrix $\boldsymbol{\Sigma}_{\text{calibrated}}$ defined in Equation \ref{eq:calibrated_covariance_matrix} assigns tight, conservative covariance values of $10^{-4}\,\text{rad}^2$ to the roll and pitch states. Since the multi-agent system operates strictly within a flat, planar $SE(2)$ manifold, keeping these out-of-plane rotational variances minimized is critical to prevent transient wheel-ground contact slips or floor seams from introducing artificial roll/pitch state drift into the real-time EKF localization server.

\section{Feature Extraction Pipelines}

\subsection{Edge Compute Trade-Offs and Preprocessing}

Executing convolutional networks or vision transformers requires substantial resources for real-time performance, which is incompatible with the power budget of a mobile robot. More specifically, a deep vision model implemented on the central processing unit (CPU) of the robot described in Section 5.1 will rapidly deplete the battery. Furthermore, the execution on the CPU may cause substantial delays for time-critical procedures like path tracking, motion control, and global localization estimators. Our system avoids this unnecessary overhead by designing a light-weight handcrafted feature extraction pipeline.

First, the system processes 1920 by 1080 color frames captured by the forward-looking camera. Each frame is down-sampled and converted to a single-channel gray-scale image according to the following luminance formula:
\[ I_{\text{gray}}(x,y) = 0.299 \cdot R(x,y) + 0.587 \cdot G(x,y) + 0.114 \cdot B(x,y) \]
This preprocessing reduces the input dimensionality while eliminating most of the high-frequency components of the original image. Next, the system extracts micro-texture and macro-structural features using two separate feature encoding pipelines.

\subsection{Spatial Rotation-Invariant Uniform Local Binary Patterns (LBP)}

To extract spatial micro-texture features from the preprocessed grayscale frames, the perception pipeline makes use of rotation-invariant uniform Local Binary Patterns. For each pixel with intensity value $g_c$, the local variation in intensities is approximated using the following local binary descriptor formulation:
\begin{equation}
LBP_{P,R}(x_c, y_c) = \sum_{p=0}^{P-1} s(g_p - g_c) \cdot 2^p \label{eq:lbp_standard}
\end{equation}
where thresholded pixel values are obtained using the binary step function:
\begin{equation}
s(z) = \begin{cases} 1 & z \ge 0 \\ 0 & z < 0 \end{cases} \label{eq:lbp_step}
\end{equation}
To reduce the descriptor's dimensionality while maintaining discriminative power, a uniformity constraint is imposed on the binary patterns. Specifically, the number of bit transitions in the binary string is calculated as:
\begin{equation}
U(LBP_{P,R}) = |s(g_{P-1} - g_c) - s(g_0 - g_c)| + \sum_{p=1}^{P-1} |s(g_p - g_c) - s(g_{p-1} - g_c)| \label{eq:lbp_uniformity}
\end{equation}
and the patterns with $U \le 2$ are selected as rotation-invariant uniform patterns. Below, a combinatoric proof shows that for a $P$-bit mask of length $P$, the rotation-invariant uniform patterns will occupy exactly $P+2$ bins after rotation-invariant grouping.

\subsubsection*{Combinatoric Proof}
Let us consider an arbitrary $P$-bit binary pattern as a circular string of bits of length $P$. Then, for each binary pattern, we calculate the number of bit transitions, i.e., the number of places where two consecutive bits are different. First, note that for $k=0$, there exists only one uniform pattern with $P$ zeros and zero bit transitions. Similarly, for $k=P$, there exists only one uniform pattern with $P$ ones and zero bit transitions.

For any sum of bits $k$ not equal to zero or $P$, it can be shown that the number of uniform patterns is equal to $P$ for each value of $k$ from $1$ to $P-1$. This follows from the fact that to form a uniform pattern of length $P$ with a sum of bits equal to $k$, we need to place $k$ consecutive 1's and $P-k$ consecutive 0's in the circular string. Since the pattern is circular, there are exactly $P$ ways to do this for each value of $k$.

Then, under rotation-invariant grouping, each rotation-invariant group of $P$ binary patterns of sum $k$ will be projected to the same bin. Therefore, there will be exactly one uniform rotation-invariant bin for each value of $k$ from $0$ to $P$, giving a total of $P+1$ bins for all rotation-invariant uniform patterns of length $P$. This is summarized in the expression:
\begin{equation}
1 + 1 + (P - 1) = P + 1 \label{eq:uniform_lbp_classes}
\end{equation}
Finally, adding one bin for the non-uniform patterns, the total number of bins for the rotation-invariant uniform patterns is equal to $P+2$, as defined by the formulation:
\begin{equation}
(P + 1) + 1 = P + 2 \label{eq:total_lbp_classes}
\end{equation}
For the value of $P=8$, the 8-bit local binary patterns are divided into exactly 10 bins under rotation-invariant grouping. 

In texture analysis formulations, the term "uniform" can sometimes denote \footnotesize{non-rotation-invariant} uniform patterns, yielding $P(P-1) + 3 = 59$ bins. However, the algorithmic implementation utilized in this design inherently incorporates circular bit-shift normalization. This operation projects all rotationally equivalent uniform patterns onto a single representative bin. Hence, the resulting feature space is mapped to the mathematically proven $P + 2 = 10$ dimensionality. This projection guarantees that the theoretical combinatoric proof aligns with the discrete histogram binning $bins = \{0, 1, \dots, P+2\}$ used to generate the final feature vector.


\subsection{Histograms of Oriented Gradients (HOG) Structural Edge Analysis}

To reveal higher-level structural features such as macro-level box deformations or structural collapse, the perception pipeline utilizes the Histograms of Oriented Gradients descriptor. The structural edge features are extracted from the preprocessed grayscale frames by calculating gradient orientations at multiple locations. Below, a dimensional analysis proves that the resulting high-dimensional vector indeed captures all of the image structure details necessary for anomaly detection.

First, consider a $128 \times 128$ image obtained after the preprocessing step. To build the HOG descriptor, the image is divided into $16 \times 16$ pixel cells, and the sliding window technique is used with $2 \times 2$ cells per block and $16$-pixel strides. The number of blocks in the horizontal and vertical directions is calculated as:
\begin{equation}
N_{\text{blocks\_x}} = \frac{W / B_{\text{cell}} - N_{\text{cell\_x}}}{S} + 1 = \frac{128/16 - 2}{1} + 1 = 7 \label{eq:hog_blocks_x}
\end{equation}
\begin{equation}
N_{\text{blocks\_y}} = \frac{H / B_{\text{cell}} - N_{\text{cell\_y}}}{S} + 1 = \frac{128/16 - 2}{1} + 1 = 7 \label{eq:hog_blocks_y}
\end{equation}
giving the total number of blocks as:
\begin{equation}
N_{\text{blocks\_total}} = N_{\text{blocks\_x}} \times N_{\text{blocks\_y}} = 7 \times 7 = 49 \label{eq:hog_blocks_total}
\end{equation}
Each block contains 4 cells, and for each cell a histogram of gradient directions is built with unsigned angles divided into 9 bins. Therefore, the dimension of the feature vector for each block is given by:
\begin{equation}
D_{\text{block}} = N_{\text{cell\_x}} \times N_{\text{cell\_y}} \times B_{\text{bins}} = 2 \times 2 \times 9 = 36 \label{eq:hog_block_dim}
\end{equation}
and the dimension of the final structural feature vector is:
\begin{equation}
D_{\text{HOG}} = N_{\text{blocks\_total}} \times D_{\text{block}} = 49 \times 36 = 1764 \label{eq:hog_total_dim}
\end{equation}
which concludes the dimensional analysis for the structural HOG features. To achieve illumination invariance, the block-wise feature vectors undergo normalization using the L2-Hys algorithm with a clipping threshold parameter of $\tau = 0.2$.

\subsection{Feature Fusion and PCA Dimensionality Reduction for Edge Detection}

To utilize these descriptors in an edge-deployed anomaly classifier, micro and macro representations can be joined together to form a fusion feature using the process of concatenation of two vectors as shown below:

\[ \mathbf{x}_{\text{fused}} = \begin{bmatrix} \mathbf{f}_{\text{LBP}}^T & \mathbf{f}_{\text{HOG}}^T \end{bmatrix}^T \in \mathbb{R}^{d_{\text{fused}}} \]

where $\mathbf{f}_{\text{LBP}} \in \mathbb{R}^{160}$ is the spatial Local Binary Pattern histogram, $\mathbf{f}_{\text{HOG}} \in \mathbb{R}^{1764}$ is the normalized gradient vector in $\text{L2-Hys}$ norm, and $d_{\text{fused}} = d_{\text{LBP}} + d_{\text{HOG}} = 1924$.

First, this combined feature vector is normalized according to the standardization algorithm using pre-fitted statistics:

\[ \mathbf{x}_{\text{scaled}, j} = \frac{\mathbf{x}_{\text{fused}, j} - \mu_j}{s_j + \epsilon} \]

where $\epsilon = 10^{-8}$ for numerical stability. Next, the system learns a linear transformation matrix according to the Principal Component Analysis (PCA) algorithm, which decreases the feature vector dimensionality, but retains the information. The linear transformation matrix $V^T \in \mathbb{R}^{1924 \times 125}$ is formed by the $k = 125$ eigenvectors associated with the largest eigenvalues of the empirical covariance matrix $\mathbf{\Sigma} \in \mathbb{R}^{1924 \times 1924}$, estimated based on the training data:

\[ \mathbf{\Sigma} = \frac{1}{M-1} \sum_{i=1}^{M} \left( \mathbf{x}_{\text{scaled}}^{(i)} \right) \left( \mathbf{x}_{\text{scaled}}^{(i)} \right)^T \]

Specifically, the $k = 125$ orthogonal vectors retain exactly 95\% of the empirical variance. The algorithm applies the linear transformation of the normalized feature vector onto the reduced eigenspace using the formula:

\[ \mathbf{x}_{\text{pca}} = \mathbf{x}_{\text{scaled}} \cdot V^T \]

Thus, the output of the PCA algorithm becomes the compressed representation of $d = 125$ dimensions. Dimensionality reduction enables the system to exclude uninformative features and multicollinearity while having a compact inference footprint. In particular, the feature vector compression rate amounts to 93.50\%. As a consequence, the system is able to continuously observe the image stream in the real-time mode with a high control loop frequency without encountering computational constraints imposed by the robot's CPU compute budget.

\subsection{Stochastic Sensor Perturbation and Environmental Noise Emulation}

A proper validation environment for the handcrafted vision pipeline is critical in evaluating its performance on the robotic edge hardware. Various nuisances inherent to the warehouse environment render the pixels in the camera's field-of-view less informative compared to ideal conditions, namely, dust on the lens, chassis vibrations, and flickering of the overhead lights. To emulate these undesirable effects, rather than applying them on the raw high resolution color frames, here we apply the three-step pipeline to the preprocessed (downsampled) 128 by 128 single-channel matrix \cite{ojala2002multiresolution}. The reason for choosing this particular resolution lies in alleviating the computational burden on the onboard computer: warping the raw 1920 by 1080 by 3 image with an affine transformation and overlaying the pixel-wise Gaussian noise would result in prohibitively large intermediate matrices in terms of both resolution and bit-depth. By shifting the computations to the 128 by 128 by 1 resolution, the memory footprint and processing time are reduced roughly by the factor of 380, which is acceptable for real-time vision processing on the embedded robot \cite{dalal2005histograms}. Essentially, the visual perception thread must not incur significant computational debt on the robotic edge hardware which has to constantly service the motor control and localization loops.

First, to emulate the nuisances of shadowing and aging of the warehouse lights, we apply a channel-wise multiplicative scaling factor:
\begin{equation} 
I_{\text{illum}}(x,y) = I_{\text{gray}}(x,y) \cdot \alpha \label{eq:noise_illum}
\end{equation}
where $\alpha$ is a scalar random variable. To allow for proper generalization and ample out-of-distribution regularization during the offline training phase, the scaling factor in Equation \ref{eq:noise_illum} is sampled from a uniform distribution spanning a large range:
\begin{equation} 
\alpha_{\text{offline}} \sim U(0.85, 1.15) \label{eq:noise_alpha_offline} 
\end{equation}
However, to avoid excessive flickering during the online tracking phase and consequent loss of consistency in visual odometry, the onboard scaling perturbation is restricted to a smaller range, capturing only the variations in ambient illumination:
\begin{equation} 
\alpha_{\text{online}} \sim U(0.95, 1.05) \label{eq:noise_alpha_online} 
\end{equation}

Next, to account for the camera tremors during the robot's emergency stops and turns, we apply a randomly rotated affine warp:
\begin{equation} 
I_{\text{jitter}} = \text{warp\_affine}(I_{\text{illum}}, \theta) \label{eq:noise_jitter} 
\end{equation}
where $\theta$ is a rotation angle sampled from a uniform distribution. To enforce strong levels of invariance to accidental sensor movements, during offline training the rotation angle $\theta$ in Equation \ref{eq:noise_jitter} is drawn from a large uniform distribution:
\begin{equation} 
\theta_{\text{offline}} \sim U(-10.0^\circ, +10.0^\circ) \label{eq:noise_theta_offline} 
\end{equation}
On the other hand, to capture only the high-frequency vibrational movements of the chassis without introducing perspective distortions during the online tracking, the rotation angle $\theta$ is sampled from a smaller range:
\begin{equation} 
\theta_{\text{online}} \sim U(-2.0^\circ, +2.0^\circ) \label{eq:noise_theta_online} 
\end{equation}

Finally, to emulate the optical sensor properties, we apply reflective padding to prevent the appearance of undesired black pixels around the image borders during the downsampling stage. To introduce stochastic sensor perturbations, we add a zero-mean white Gaussian noise with a standard deviation $\sigma_{\text{noise}}$:
\begin{equation} 
I_{\text{sensor}}(x,y) = I_{\text{jitter}}(x,y) + \mathcal{N}(0, \sigma_{\text{noise}}^2) \label{eq:noise_sensor} 
\end{equation}
where the noise standard deviation is dynamically adjusted according to:
\begin{equation} 
\sigma_{\text{noise}} \sim U(0.0, 5.0) \label{eq:noise_sigma} 
\end{equation}
and the resulting intensity values are clipped to the 0-255 range. This final transformation is critical in endowing the vision module with the ability to deal with extreme lighting conditions, and thus it is applied right before the LBP descriptor is extracted from the image \cite{mei2005case}.

\subsection{Computational Complexity and Latency Analysis}

In order to formally establish the real-time capability of the proposed handcrafted perception pipeline on-board a low-power embedded processor, in this section we carry out an asymptotic complexity analysis of the involved operations. In particular, we characterize the time scalability of each stage by expressing their computational cost in terms of Big-O notation.

Assume that the image resolution captured by the optical sensor is given by $W_{\text{raw}} \times H_{\text{raw}}$ while the downsampled single-channel resolution is $W \times H$. First, the grayscale preprocessing stage has linear complexity in terms of input-output pixels, given by $\mathcal{O}(W_{\text{raw}} \cdot H_{\text{raw}} + W \cdot H)$. Since the spatial uniform LBP descriptor requires to estimate a vector of size $P$ for each pixel in the downsampled image, while mapping and pooling such statistics in the localized spatial histograms, its overall complexity is also linear in space, scaling as $\mathcal{O}(W \cdot H \cdot P)$.

On the other hand, HOG feature extraction involves gradient calculations and pooling statistics of size $B_{\text{bins}}$ in blocks, resulting in a linear complexity in space that scales as $\mathcal{O}(W \cdot H \cdot B_{\text{bins}})$. The following feature-level concatenation and standardization steps have linear complexity of $\mathcal{O}(d_{\text{fused}})$ in terms of the concatenated feature dimension, where $d_{\text{fused}} = 1924$. 

Projection on the reduced PCA space follows directly from the matrix-vector multiplication between the $d_{\text{fused}} \times d_{\text{pca}}$ weight matrix and the standardized feature vector, also with a linear complexity scaling as $\mathcal{O}(d_{\text{fused}} \cdot d_{\text{pca}})$, where $d_{\text{pca}} = 125$. Finally, the linear SVM classification amounts to estimating the inner product between the PCA projected feature vector and the hyper-plane weight vector, which again has a linear complexity of $\mathcal{O}(d_{\text{pca}})$ in terms of $d_{\text{pca}}$.

Combining these individual bounds, the total asymptotic complexity of the complete onboard perception pipeline is given by:
\begin{equation}
\mathcal{O}(W \cdot H \cdot P + d_{\text{fused}} \cdot d_{\text{pca}}) \label{eq:total_pipeline_complexity}
\end{equation}
Thanks to the low-resolution input ($128 \times 128$), a small number of neighborhood pixels ($P = 8$), and reduced feature dimension ($d_{\text{fused}} = 1924$, $d_{\text{pca}} = 125$), the complexity of each individual stage is kept well within acceptable limits. In particular, we find that the total execution time of the \texttt{waypoint\_sender.py} node is less than 1.2 ms with no significant CPU-time overhead, thus enabling the parallel execution of path-tracking and localization routines.

\section{Edge Classifier Design and Custom Spatial Back-Projection XAI}

\subsection{Linear Support Vector Machine Classification}

The cardboard anomaly classification task on the edge is particularly challenging due to the need to reach a high-confidence decision at the fastest possible time. To this end, the reduced feature vector representation is scored using a soft-margin Linear Support Vector Machine. Instead of directly optimizing the primal problem, which for the case of high-dimensional input bounds can scale poorly, the optimal separating hyperplane is resolved in the dual formulation as:
\begin{equation}
\max_{\boldsymbol{\alpha}} \sum_{i=1}^{M} \alpha_i - \frac{1}{2} \sum_{i=1}^{M} \sum_{j=1}^{M} \alpha_i \alpha_j y_i y_j \langle \mathbf{x}_i, \mathbf{x}_j \rangle \label{eq:svm_lagrangian_dual}
\end{equation}
\text{subject to:}
\begin{equation}
\sum_{i=1}^{M} \alpha_i y_i = 0, \quad 0 \le \alpha_i \le C_{\text{svm}} \label{eq:svm_constraints}
\end{equation}
where the Lagrange multipliers are expressed as $\alpha_i$, the anomaly targets take the form of $y_i \in \{-1, \; +1\}$, and the input PCA-reduced feature vector is given by $\mathbf{x}_i \in \mathbb{R}^{125}$. The soft-margin regularization penalty, which introduces the trade-off between the model's training error and margin maximization, is controlled by a scalar hyperparameter $C_{\text{svm}} = 0.1$. After obtaining the optimal Lagrangian dual variables, the decision hyperplane parameters are recovered, and the continuous decision score for an arbitrary online input feature vector $F_{\text{pca}}$ is computed as:
\begin{equation}
f(F_{\text{pca}}) = \mathbf{w}_{\text{pca}} \cdot F_{\text{pca}} + b \label{eq:svm_decision_function}
\end{equation}
with the bias term $b$. If the decision value $f(F_{\text{pca}})$ is positive, the target cardboard packaging is considered damaged; otherwise, it is classified as intact.

\subsection{Mathematical Back-Projection to Raw Featurespace}

While being effective in automated cardboard packaging inspection, such an approach lacks the ability to offer visual explanations of the decisions to the human supervisor, thus potentially undermining the overall cyber-physical system's trustworthiness. The common model-agnostic post-hoc explanation techniques, such as input perturbation, require an unacceptable amount of computational resources for the embedded robotic vision application. To resolve the conundrum, the interpretability is baked into the model by leveraging the linearity of the reduced-space classifier's decision boundary.

Due to the SVM decision function's linearity in the input features reduced by PCA, the optimal separating hyperplane's weight vector can be back-projected to the original high-dimensional feature space using the same projection matrix utilized for the forward transformation. Specifically, the raw feature-space equivalent of the reduced-space weight vector $\mathbf{w}_{\text{pca}}$ is calculated as:
\begin{equation}
\mathbf{w}_{\text{raw}} = \mathbf{w}_{\text{pca}} \mathbf{V} \label{eq:xai_back_projection}
\end{equation}
To account for the standardization step, the raw weight coefficients are scaled by the inverse of the standard deviations:
\begin{equation}
w_{\text{scaled}, j} = \frac{w_{\text{raw}, j}}{s_j + \epsilon} \label{eq:scaling_reversal}
\end{equation}
where epsilon is set to $\epsilon = 10^{-8}$ for numerical stability. This allows recovering the attribution weights that reflect the contribution of each handcrafted feature to the final decision.

\subsection{Spatial Attributions and Image Blending}

To obtain the spatial attribution weights relevant to the local texture statistics of the cardboard box surface, the first 160 elements of the scaled weight vector $\mathbf{w}_{\text{scaled}}$ corresponding to the uniform LBP block features are extracted, while the remaining 1764 dimensions representing the HOG edge features are intentionally bypassed. 

This exclusion is mathematically mandated by the spatial properties of the two descriptors: LBP features are extracted from strict, non-overlapping spatial grid blocks ($4 \times 4$ regions), allowing an isomorphic, one-to-one mapping back to the 2D image coordinate plane. Conversely, HOG features capture normalized gradient orientations ($\theta_{\text{gradient}}$) pooled over heavily overlapping block regions with spatial normalization. Attempting to back-project these normalized gradient histograms directly to raw 2D pixel coordinates would induce severe spatial blur, phase misalignment, and destructive gradient interference, rendering the visual explanation map uninterpretable. Thus, isolating LBP magnitudes ensures a geometrically sound and localized spatial attention map without loss of explainable fidelity. After thresholding, the absolute values are concatenated into a weight vector:
\begin{equation}
\mathbf{w}_{\text{LBP}} = |\mathbf{w}_{\text{scaled}}[0 : 160]| \label{eq:lbp_weights_isolation}
\end{equation}
where the block contributions are summed over the channel dimension:
\begin{equation}
I_{\text{block}, k} = \sum_{m=1}^{10} \mathbf{w}_{\text{LBP}}[(k \cdot 10) + m] \label{eq:lbp_block_sum}
\end{equation}
resulting in a 16-element vector reshaped into a $4 \times 4$ matrix $\mathbf{M}_{\text{grid}} \in \mathbb{R}^{4 \times 4}$ interpolated to the original 128 by 128 image resolution with a bilinear kernel to form a jet colormap attention map. The attention map is then overlaid on the original sensor preprocessing frame, which is converted to a three-channel grayscale version:
\begin{equation}
I_{\text{gray, BGR}}(x,y) = \begin{bmatrix} I_{\text{gray}}(x,y) & I_{\text{gray}}(x,y) & I_{\text{gray}}(x,y) \end{bmatrix}^T \in \mathbb{R}^3 \label{eq:grayscale_expansion}
\end{equation}
in a linear blending operation that produces the final XAI visualization:
\begin{equation}
I_{\text{XAI}}(x,y) = (1 - \alpha_{\text{blend}}) \cdot I_{\text{gray, BGR}}(x,y) + \alpha_{\text{blend}} \cdot I_{\text{JetMap}}(x,y) \label{eq:alpha_blending}
\end{equation}
where the blending coefficient alpha is set to $\alpha_{\text{blend}} = 0.35$. Notably, the full feature-space attribution pipeline runs on the robotic edge vision system in only 1.2 ms, thus enabling real-time attention-aware robotic box inspection.

\subsection{Steady-State Foundation for Resilient Fleet Coordination}

The mathematical, geometric, and computational steady-state models introduced in this chapter-governing battery consumption, spatial mutual exclusion, parametric waypoint calculation, non-holonomic tracking, and edge perception-collectively constitute the nominal operating envelope of the fleet. Under this framework, any physical or logical deviation from these baseline steady-state values can be utilized as a deterministic failure detection mechanism. Consequently, this steady-state operating envelope serves as the foundation for developing the hierarchical fault-tolerant control policies, complex recovery decision-making logic, and decentralized multi-robot task reallocation algorithms detailed in Chapter 6. These couplings emerge through specific mathematical relationships. Namely, the hybrid battery consumption model formulated in Equation~\ref{eq:hybrid_decay_discrete} supplies the critical telemetry input to the battery-proportional decentralized task reallocation algorithm. In a similar vein, the linear support vector machine classification dual formulation stated in Equation~\ref{eq:svm_lagrangian_dual} provides the baseline for the decision threshold. A prolonged deviation or downgrading of classification performance at this level triggers the analytical redundancy fallback mechanism, thereby initiating a failover to the optimized XGBoost classifier model. In this regard, it is possible to state that the steady-state formulations supply the activation thresholds to the resilient control loops.